\documentclass[a4paper,10pt,lithuanian]{article}

\usepackage{lnfkstyle}

\newcommand{\e}[1]{\mathbf{e}_{#1}}
\newcommand{\sig}[1]{\boldsymbol{\sigma}_{#1}}
\newcommand{\Isig}[1]{I\!\boldsymbol{\sigma}_{#1}}
\newcommand{\ii}{\mathrm{i}} %imaginary unit
\newcommand{\ee}{\mathrm{e}} %exponent
\def\eps#1{\varepsilon_{#1}} %dielectric constants

%Cliffordo algebros simboliai
\def\d{\!\cdot\!} %dot product
\def\D{\cdot}
\def\w{\!\wedge\!} %wedge product
\def\W{\wedge}
\def\e#1{\mathbf{e}_{#1}} %Clifford basis. Input: $\e3$ or $\e{12}$
\def\g#1{\gamma_{#1}} %CL31 basis. Input: $\g3$ or $\g{12}$
\def\Ig#1{I\!\gamma_{#1}}%CL31 trivectors
\def\sig#1{\boldsymbol{\sigma}_{#1}} %CL30 basis vectors. $\sig{12}$
\def\Isig#1{I\!\boldsymbol{\sigma}_{#1}}


\def\bE{\mathbf{E}}
\def\bH{\mathbf{H}}
\def\bD{\mathbf{D}}
\def\bF{\mathbf{F}}
\def\bB{\mathbf{B}}
\def\bW{\mathbf{W}}

\def\ba{\mathbf{a}}
\def\bb{\mathbf{b}}
\def\bk{\mathbf{k}}
\def\bx{\mathbf{x}}
\def\bw{\mathbf{w}}

\def\cF{\mathcal{F}}
\def\cD{\mathcal{D}}
\def\cE{\mathcal{E}}
\def\cG{\mathcal{G}}
\def\cB{\mathcal{B}}
\def\cH{\mathcal{H}}
\def\cW{\mathcal{W}}

\def\beps{\boldsymbol{\varepsilon}}
\def\bmu{\boldsymbol{\mu}}



\begin{document}

\titlelt{Medžiaginiai ryšiai klasikinėje ir reliatyvistinėje optikoje GA požiūriu}

\titleen{Classical and relativistic constitutive relations from GA point of view}


%\author{\uline{Adolfas Dargys}$^{1}$}
\author{Adolfas Dargys}

\keywords{medžiaginiai sąryšiai, elektromagnetinių bangų sklidimas,  Cliffordo algebra}

\maketitle

\address{$^{1}$Nacionalinis fizinių ir technologijos mokslų centras,  Puslaidininkių fizikos institutas, Saulėtekio 3, LT-10257 Vilnius}

%\address{$^{2}$Fizinių ir technologijos mokslų centras, Savanorių pr. 231, 02300 Vilnius}

\rightaddress{adolfas.dargys@ftmc.lt}

\begin{multicols*}{2}
Diferencialinių Maxwell'o  lygčių sistema nėra pilna.  Ją reikia
papildyti medžiaginiais ryšiais tarp pirminių laukų $(\bE,\bB$),
elektrinio ir magnetinio, ir medžiagoje indukuotų antrinių laukų
$(\bD,\bH)$~\cite{key-1}. Ryšiai gali nusakyti medžiagos inertiškumą ar
netiesiškumą ją žadinant, histerezę ir pan. Paprasčiausiu atveju
manoma, kad  ryšis tarp pirminių laukų $(\bE,\bB$) ir  indukuotų medžiagoje $(\bD,\bH)$ yra
momentinis. Klasikinės elektrodinamikos atveju medžiaginiai ryšiai
tokiais atvejais užrašomi per $3\times 3$ matricas
$\bar{\varepsilon}$, $\bar{\gamma}$, $\bar{\beta}$ ir
$\bar{\mu}^{-1}$~\cite{key-1}:
\begin{equation}\begin{split}
&\bD=\bar{\varepsilon}\bE+\bar{\gamma}\bB,\\
&\bH=\bar{\beta}\bE+\bar{\mu}^{-1}\bB.
 \end{split}\end{equation}
Kyla klausimas, koks  bendriausias tokių matricų pavidalas ir kaip jas konstruoti  tuo atveju,
kai medžiaga yra homogeniška. Į tokius klausimus leidžia atsakyti
Cliffordo (\textit{Cl}) geometrinė algebra (GA): klasikiniu atveju
\textit{Cl}$_{3,0}$, o reliatyvistiniu atveju \textit{Cl}$_{1,3}$,
kur apatiniai indeksai nusako erdvės (medžiagos) metriką~\cite{key-2}. Mūsų
darbuose \cite{key-2,key-3} parodyta, kad elektromagnetinių bangų sklidimo
terpėse  savybės tokiais atvejais seka iš algebrų
\textit{Cl}$_{3,0}$ ir \textit{Cl}$_{1,3}$  vidinės sandaros, ir
todėl medžiaginiai ryšiai, t.~y. matricų $\bar{\varepsilon}$,
$\bar{\gamma}$, $\bar{\beta}$ ir $\bar{\mu}^{-1}$ pavidalas, seka
iš GA multivektorių simetrijos savybių.  Konkrečiai PT simetrija įskaityta \textit{Cl}$_{3,0}$,
o erdvėlaikio CPT simetrija  --  \textit{Cl}$_{1,3}$  algebroje.   Darbe~\cite{key-3} pateikti bendri
medžiagų saryšiai klasikinės elektrodinamikos atvejui, t.~y. kai
terpės greitis yra žymiai mažesnis už šviesos greitį, o
straipsniuose \cite{key-4,key-5} -- reliatyvistinei elektrodinamikai.

Reliatyvistinės eletrodinamikos atveju turime vienintelį lauką,
taip vadinamą Faraday'aus lauką medžiagoje $\cF=\cE+\cB$.
Geometrinėje algebroje laukai nusakomi bivektoriais (orientuotomis
plokštumomis). Plokštumos dydis nusako lauko stiprį, o plokštumos
orientacija -- lauko kryptį. Matuojami elektrinis $\cE$  ir  magnetinis $\cB$ laukai
priklauso nuo stebėtojo. Invariantu išlieka tik ju suma  $\cF$.
Panašiai užrašomas sužadintas  medžiagoje laukas,
$\cG=\cD+\cH=\chi(\cF)$, kur $\chi$ yra apibendrinta medžiagos
juta, kurią galima susieti su medžiaginėmis matricomis $\bar{\varepsilon}$, $\bar{\gamma}$, $\bar{\beta}$ ir
$\bar{\mu}^{-1}$ formulėje~(1). Bivektoriai $\cF$ ir $\cG$ turi po šešias
dedamąsias, todėl medžiagos sąryšius yra patogu nusakyti $6\times
6$ matricomos:
\begin{equation}\begin{split}\label{transform}
 \left[\begin{array}{c}
\cD\\
\cH\\
\end{array}\right]&=\Big(
\left|\begin{array}{c|c}
\text{Diel.Birefr.}&\text{Fizeau}\\
\hline
\text{Fizeau}&\text{Magn.Birefr.}\\
\end{array}\right|
+\\
%%
&\quad\left|\begin{array}{c|c}
\text{Diel.Faraday}&\text{Opt. Aktyv.}\\
\hline
\text{Opt. Aktyv.}&\text{Magn.Faraday}\\
\end{array}\right|
\Big) \left[\begin{array}{c}
\cE\\
\cB\\
\end{array}\right|.
\end{split}\end{equation}
Į kiekvieną iš $3\times 3$ blokų įrašėmę po vieną jam charakteringą reiškinį
(Fizeau $\rightarrow$ Fizeau reiškinys, Birefr. $\rightarrow$
dvejopas lūžimas, Opt. Aktyv. $\rightarrow$ optinis aktyvumas). Bivektoriai $\cE$ ir $\cD$ yra erdviškieji, o
bivektoriai $\cB$ ir $\cH$ -- laikiškieji, t.~y. jų kvadratai yra
atitinkamai teigiami/neigiami dydžiai. Taigi $6\times 6$ matricą,
kurios pirmoji dalis yra simetrinė, o antroji antisimetrinė,
galima interpretuoti kaip bendrojo bivektoriaus $\cF$
transformaciją į kitą bivektorių. Kadangi \textit{Cl}$_{1,3}$
algebra automatiškai įskaito erdvėlaikio simetrijos savybes (CPT
simetriją) per involiucijas, atskiri matricų blokai (2) formulėje
taip pat pasižymi tam tikromis simetrijos  savybėmis, kurios ir
nusako mežiagos reliatyvistinius sąryšius. Atitinkamos $6\times 6$
matricos (jų yra penkios) yra suskaičiuotos su geometrine Cliffordo algebra  ir   pateiktos straipsniuose~\cite{key-4,key-5},
kurias sudedant įvairiomis kombinacijomis galima gauti pačius
įvairiausius medžiaginius sąryšius, kurie seka iš relatyvistinės
elektrodinamikos. Iš pateiktų visų galimų reliatyvistinių $6\times
6$ matricų galima sukonstruoti  klasikinio pavidalo  (1) sąryšius, kurie
dažniausiai ir nauduojami bangų sklidimo analizėje. Smulkiai tokie
sąryšiai išrašyti~\cite{key-5} darbe. Juos galima palyginti su klasikinės elektrodinamikos ryšiais~\cite{key-3},
iš kur seka,  kad klasikinė elektrodinamika įskaito ne visus bangų sklidimo reiškinius.

  
  
  

\begin{thebibliography}{References}

\bibitem{key-1}M.~Born and E.~Wolf, 1999 \textit{Principles of Optics}
   (Cambridge, CUP), 1999.

\bibitem{key-2}A.~Dargys, A.~Acus, 2015 \textit{Cliffordo
  geometrinė algebra  ir jos taikymai} (UAB Petro ofsetas), 2015.

\bibitem{key-3}A.~Dargys,   Lith. J. Phys. {\bf 55}, 92-99 (2015)

\bibitem{key-4}A.~Dargys, Opt. Comm. {\bf 354}, 259-65 (2015)

\bibitem{key-5}A.~Dargys,  arXiv:1609:04261v1, [physics.class-ph] (2016)

\end{thebibliography}





\end{multicols*}
\end{document}
